% Section 4.1, from lectures/L04/appendix/a_modules.md.

\section{Kernel Modules, and the C You Are Allowed to Write}
\label{c4:app:a}

This section is the whole of the chapter's material: what a module is, how it is built and loaded,
what the kernel does with the licence you declare, and what the C environment inside a kernel is
actually like. It ends with the coding style question, which this repository answers differently
from the kernel and owes you an explanation for.

\subsection{A module is not a program}
\label{c4:app:a1}

The smallest useful module is six lines of substance:

\begin{ccode}
// SPDX-License-Identifier: GPL-2.0-only
#include <linux/init.h>
#include <linux/module.h>

static int __init qa_hello_init(void)
{
    pr_info("qa_hello: loaded\n");
    return 0;
}

static void __exit qa_hello_exit(void)
{
    pr_info("qa_hello: unloaded\n");
}

module_init(qa_hello_init);
module_exit(qa_hello_exit);
MODULE_LICENSE("GPL");
\end{ccode}

\noindent There is no \code{main}, and \code{qa_hello_init} is not one. The distinction is worth
being precise about because it shapes everything you write afterwards.

\code{main} runs, and while it runs, your program exists. When it returns, your program is over.
\code{qa_hello_init} runs \textbf{once, at load}, and when it returns the module is \textbf{resident
and idle}. Nothing of yours is running. The module is a body of code and data sitting in the
kernel's address space, waiting to be called by something else: a file operation, an interrupt, a
timer, a probe.

So the shape of a module is \emph{registration}. \code{init} announces what you can do and then gets
out of the way; everything after that is callbacks. A module whose \code{init} function contains a
loop is almost always a module that should have started a thread, and a module whose \code{init}
never returns has hung the process that called \code{insmod}.

\textbf{\code{__init} and \code{__exit} are instructions to the linker.} Code marked \code{__init}
is placed in a section that is freed once initialisation is done, which is why you see \code{Freeing
unused kernel memory: 4672K} at the end of boot. \code{__exit} code can never run when the module is
built into the kernel rather than as a module, because a built-in module can never be removed, so it
is thrown away: at link time on most architectures, and on arm64 at the end of boot, with the
\code{__init} code.

\lecturefigure{lectures/L04/appendix/images/module_lifecycle.png}
  {Five boxes in a row: insmod, init, resident, exit and rmmod. The middle box is highlighted and
   annotated to say that nothing of yours runs there; the module is code and data occupying kernel
   memory, waiting to be called by a system call, an interrupt, a timer or another module.}
  {fig:module_lifecycle}

\subsection{The four macros}
\label{c4:app:a2}

\begin{booktable}{@{}>{\hsize=0.43\hsize}Ll>{\hsize=1.57\hsize}L@{}}
  \thead{Macro} & \thead{Required} & \thead{What it is for} \\
  \midrule
  \code{MODULE_LICENSE} & \textbf{Yes} & Read by the loader. See \secref{c4:app:a5}. Its absence is an error at build time \\
  \code{MODULE_AUTHOR} & No & Convention, and useful when a bug report reaches a maintainer \\
  \code{MODULE_DESCRIPTION} & No & Shown by \code{modinfo} \\
  \code{MODULE_VERSION} & No & Free text; not used for dependency resolution \\
\end{booktable}

\noindent \code{MODULE_LICENSE} is the one that is not optional. Omit it and \code{modpost} fails
the build with \code{missing MODULE_LICENSE()}, which is a good error, because the alternative is a
module that loads as if it were proprietary: it taints the kernel and cannot see a single GPL-only
symbol.

All of them are readable without loading anything:

\begin{console}
$ modinfo qa_hello.ko
filename:       qa_hello.ko
license:        GPL
depends:
vermagic:       6.12.30 SMP preempt mod_unload modversions aarch64
\end{console}

\noindent That \code{vermagic} line is \secref{c4:app:a3}.

\subsection{Building out of tree, and why a \code{.ko} is not portable}
\label{c4:app:a3}

An out-of-tree module is built by running the \emph{kernel's} build system on your directory:

\begin{shell}
make -C /path/to/kernel M=/path/to/your/module modules
\end{shell}

\noindent You are not compiling against kernel headers with your own build system. You are asking
kbuild to compile your file with exactly the flags the kernel was compiled with, which is why the
kernel tree has to be configured and built first, and why there is no shorter way.

The consequence is \code{vermagic}. Every module records the kernel version, whether that kernel was
SMP, which preemption model it used, and the architecture. The loader checks all of it and refuses a
mismatch. This is a module built against this course's \file{build/kernel}, loaded into the
real-time kernel Chapter~\ref{c12:ch} builds from the same source:

\begin{console}
qa_hello: version magic '6.12.30 SMP preempt mod_unload modversions aarch64' should be
          '6.12.30 SMP preempt_rt mod_unload modversions aarch64'
\end{console}

\noindent This looks like pedantry and is not. The kernel has no stable internal ABI; a structure
that gained a field between two point releases would be read at the wrong offsets by a module
compiled against the old one, and the failure would be silent corruption rather than a refusal to
load. \code{vermagic} turns that into an error message.

With \code{CONFIG_MODVERSIONS}, which this course enables, the check is finer: a checksum is
computed per exported symbol from its argument types, a module is refused if any symbol it uses has
changed shape, and the release number at the front of \code{vermagic} is no longer compared at all.
That is what lets a distribution ship third-party modules across a stable series.

\textbf{This course's Kbuild files build every \code{.c} in the lab directory as its own module:}

\begin{makecode}
obj-m += $(patsubst $(src)/%.c,%.o,$(filter-out %.mod.c,$(wildcard $(src)/*.c)))
\end{makecode}

\noindent That is unusual for real kernel code, where the object list is written out explicitly. It
is done here because the repository ships the build file and you write the sources, so a fixed list
would mean editing the build every time you start an exercise. Note \code{$(src)} rather than a bare
wildcard: kbuild runs the file from the kernel's output directory, and a wildcard over the current
directory silently matches nothing, producing a build that succeeds having compiled zero modules.
The \code{filter-out} is the other half: kbuild generates a \file{foo.mod.c} beside every
\file{foo.c} it builds, and without it the second build tries to make a module of that as well.

\subsection{Four commands, and the two people confuse}
\label{c4:app:a4}

\begin{booktable}{@{}llL@{}}
  \thead{Command} & \thead{Takes} & \thead{Does} \\
  \midrule
  \code{insmod} & A path & Loads exactly that file. Resolves nothing \\
  \code{modprobe} & A name & Looks the name up in \code{modules.dep}, loads dependencies first, then it \\
  \code{depmod} & Nothing & Scans \code{/lib/modules/$(uname -r)} and writes \code{modules.dep} \\
  \code{rmmod} & A name & Unloads, if nothing holds a reference \\
\end{booktable}

\noindent The demonstration that makes the difference concrete is to load a module with a dependency
the wrong way:

\begin{console}
# insmod ./qa_user.ko
qa_user: Unknown symbol qa_answer (err -2)
insmod: can't insert './qa_user.ko': unknown symbol in module or invalid parameter
\end{console}

\noindent \code{insmod} did what it was told. It loaded that file, found an unresolved symbol, and
gave up. It never occurred to it to look for a module providing \code{qa_answer}, because looking
things up is \code{modprobe}'s job and \code{modprobe} reads a file that \code{depmod} generated.

\textbf{Unloading is refcounted.} While \code{qa_user} is loaded, \code{qa_export} cannot be
removed:

\begin{console}
# rmmod qa_export
rmmod: remove 'qa_export': Resource temporarily unavailable
# cat /sys/module/qa_export/refcnt
1
# ls /sys/module/qa_export/holders/
qa_user
\end{console}

\noindent The \code{holders} directory is the useful one: it names who is holding you, rather than
merely how many. Note also that \textbf{BusyBox's \code{rmmod} returns success even when the removal
failed}, which is why the lab for this chapter checks whether the module is still loaded rather
than checking an exit status. A test that trusts that exit status gets the same answer on a kernel
that refused and on a kernel that did not.

\subsection{The licence, and where the check happens}
\label{c4:app:a5}

\Secref{c1:app:b4} made the argument. This is the mechanism.

Symbols leave a module in one of two ways:

\begin{ccode}
EXPORT_SYMBOL(qa_answer);       /* any module */
EXPORT_SYMBOL_GPL(qa_answer);   /* GPL-compatible modules only */
\end{ccode}

\noindent The accepted values of \code{MODULE_LICENSE} are a fixed list: \code{"GPL"}, \code{"GPL
v2"}, \code{"GPL and additional rights"}, \code{"Dual BSD/GPL"}, \code{"Dual MIT/GPL"}, \code{"Dual
MPL/GPL"}, and \code{"Proprietary"}. Anything else is treated as proprietary, including a typo.

Now the part worth doing rather than reading. Take a module that calls a \code{_GPL} symbol, declare
its licence \code{"Proprietary"}, and build it. \textbf{The build fails}, and the message is
unusually good:

\begin{console}
ERROR: modpost: GPL-incompatible module qa_proprietary.ko uses GPL-only symbol 'qa_answer'
\end{console}

\noindent That is \code{modpost}, a build step that runs after every module is linked and before the
\code{.ko} is finished. It reads the exports of everything it can see, and the check is three lines
in \file{scripts/mod/modpost.c}:

\begin{ccode}
if (!mod->is_gpl_compatible && exp->is_gpl_only)
        error("GPL-incompatible module %s.ko uses GPL-only symbol '%s'\n",
              basename, exp->name);
\end{ccode}

\noindent \textbf{The licence is enforced twice, at two different stages, and it is worth knowing
both} because they fail differently.

\textbf{At build time}, as above, when \code{modpost} can see that the symbol is GPL-only. That is
the usual case: the exports of the built-in kernel come from its \code{Module.symvers}, and modules
built together see each other. You get a clear message naming the licence, the module and the
symbol, and no \code{.ko} is produced at all.

\textbf{At load time}, as a backstop, and here the failure is much less obvious. The module loader
does not check the licence of a symbol you asked for and refuse it; instead it \textbf{hides
GPL-only symbols from the search entirely}. From \file{kernel/module/main.c}:

\begin{ccode}
static bool find_exported_symbol_in_section(...)
{
        if (!fsa->gplok && syms->license == GPL_ONLY)
                return false;
\end{ccode}

\noindent where \code{gplok} is set from the loading module's own taint:

\begin{ccode}
.gplok = !(mod->taints & (1 << TAINT_PROPRIETARY_MODULE)),
\end{ccode}

\noindent So a proprietary module looking up a GPL-only symbol is told the symbol \textbf{does not
exist}, and \code{insmod} reports:

\begin{console}
qa_user: Unknown symbol qa_answer (err -2)
\end{console}

\noindent \textbf{Nothing in that message mentions a licence.} It is the same message you get for a
genuine typo or a missing dependency, which makes it one of the more confusing failures in kernel
work. The load-time path is the one you hit when \code{modpost} could not have known: a module built
with \code{KBUILD_MODPOST_WARN=1}, which turns an unresolved symbol from an error into a warning, or
a module built against a kernel where the symbol was \code{EXPORT_SYMBOL} and loaded into one where
it has since become \code{EXPORT_SYMBOL_GPL}. That second case is real and happens on kernel
upgrades.

\textbf{Tainting is separate and additive.} Loading any out-of-tree module taints the kernel, even a
GPL one:

\begin{console}
qa_hello: loading out-of-tree module taints kernel.
\end{console}

\noindent \file{/proc/sys/kernel/tainted} is a bitmask;
\file{Documentation/admin-guide/tainted-kernels.rst} decodes it. Loading an out-of-tree module sets
bit 12 (\code{TAINT_OOT_MODULE}, value 4096); loading a proprietary one \emph{also} sets bit 0
(\code{TAINT_PROPRIETARY_MODULE}), so a proprietary out-of-tree module leaves
\file{/proc/sys/kernel/tainted} reading \textbf{4097} and shows as \code{(PO)} in
\file{/proc/modules}.

What a taint means is mostly social: an oops from a tainted kernel is one that upstream maintainers
will decline to debug, because they cannot see the code that might have caused it.

\textbf{One taint has a concrete technical consequence, and it is worth knowing before
Chapter~\ref{c7:ch}.} Loading a proprietary module prints this as well:

\begin{console}
qa_prop_ok: module license 'Proprietary' taints kernel.
Disabling lock debugging due to kernel taint
\end{console}

\noindent The lock validator is switched off, for the rest of that boot, because it cannot reason
about locking in code it cannot see. Chapter~\ref{c7:ch} is built entirely on that validator. So a
single proprietary module loaded anywhere on the system silently removes the tool you would use to
find the deadlock it might be causing, and the rule that follows is to keep development kernels free
of proprietary modules even when the production system has one.

\subsection{Parameters, and sysfs you did not write}
\label{c4:app:a6}

One line:

\begin{ccode}
static char *who = "world";
module_param(who, charp, 0444);
MODULE_PARM_DESC(who, "who to greet");
\end{ccode}

\noindent buys a load-time argument and a sysfs file:

\begin{console}
# insmod ./qa_hello.ko who=embedded
# cat /sys/module/qa_hello/parameters/who
embedded
\end{console}

\noindent Nobody wrote code to create that file. The third argument is its mode, and \code{0} means
do not create one at all. If the mode is writable, userspace can change the variable while the
module is loaded, which is worth pausing on: that is a shared variable being written by another
context, and everything Chapter~\ref{c7:ch} says about concurrency applies to it.

This is the first appearance of the pattern the whole driver model rests on, which is that you
describe something and the kernel generates the interface. Chapter~\ref{c10:ch} is where it stops
looking like a convenience and starts looking like the architecture.

\subsection{Printing}
\label{c4:app:a7}

\code{printk} with a level, or the wrappers, which is what you should actually use:

\begin{booktable}{@{}lrL@{}}
  \thead{Wrapper} & \thead{Level} & \thead{Use for} \\
  \midrule
  \code{pr_emerg} & 0 & The system is unusable \\
  \code{pr_alert} & 1 & Action must be taken immediately \\
  \code{pr_crit} & 2 & Critical conditions \\
  \code{pr_err} & 3 & An error your driver could not handle \\
  \code{pr_warn} & 4 & Something unexpected that you did handle \\
  \code{pr_notice} & 5 & Normal but significant \\
  \code{pr_info} & 6 & Normal. Most of what you write during development \\
  \code{pr_debug} & 7 & Compiled out unless dynamic debug is on \\
\end{booktable}

\noindent Three practical notes.

\textbf{Set \code{pr_fmt} once, at the top of the file}, and every message gets a prefix for free:

\begin{ccode}
#define pr_fmt(fmt) KBUILD_MODNAME ": " fmt
\end{ccode}

\noindent \textbf{\code{pr_debug} is not \code{pr_info} you comment out.} With
\code{CONFIG_DYNAMIC_DEBUG}, which this course enables, each call site can be switched on at run
time:

\begin{shell}
echo 'module qa_hello +p' > /sys/kernel/debug/dynamic_debug/control
\end{shell}

\noindent Until then it costs nothing but a few bytes of metadata. That is what makes it reasonable
to leave debug printing in shipped code.

\textbf{Prefer \code{dev_info} once you have a \code{struct device}}, from Chapter~\ref{c10:ch}
onwards. It prints which device the message came from, which matters the moment there are two of
them.

\subsection{The C you are allowed to write}
\label{c4:app:a8}

The kernel is compiled \code{-nostdinc}, \code{-ffreestanding}, against its own headers. There is no
C library. The practical list:
\begin{itemize}
  \item \textbf{No \file{stdio.h}, \file{stdlib.h}, \file{string.h}.} There are kernel equivalents
    with the same names for some things (\code{memcpy}, \code{strlen}, \code{snprintf}) and no
    equivalent for others.
  \item \textbf{No \code{malloc}.} \code{kmalloc} and friends, which take a flag saying what context
    you are in. Chapter~\ref{c6:ch}.
  \item \textbf{No floating point.} The kernel does not save FPU state on entry, so using a
    \code{double} corrupts whatever userspace task was interrupted. On arm64 the compiler is not
    allowed to try: the kernel is built with \code{-mgeneral-regs-only}, so a \code{double} is a
    compile error. If you genuinely need it there is \code{kernel_fpu_begin()}, and needing it is
    nearly always a sign the calculation belongs in userspace. Fixed point is the normal answer.
  \item \textbf{A small, fixed stack.} 16 KB on arm64, for the whole call chain, and nothing grows
    it; interrupts run on a separate per-CPU stack of the same size. This kernel puts an unmapped
    guard page below each stack (\code{CONFIG_VMAP_STACK}, on by default), so running off the end is
    a \code{kernel stack overflow} panic rather than a diagnostic you can recover from. Without the
    guard page, a large local array silently corrupts whatever lies below the stack, and the crash
    comes somewhere else entirely. \code{-Wframe-larger-than=2048} warns about the obvious cases at
    compile time.
  \item \textbf{No exceptions and no unwinding.} Every function returns an error code, and the
    caller checks it. The convention is zero for success and a negative \code{-Exxx} for failure,
    and pointer-returning functions use \code{ERR_PTR}, \code{IS_ERR} and \code{PTR_ERR} to pack an
    error into a pointer.
  \item \textbf{Never trust a userspace pointer.} Chapter~\ref{c5:ch}.
\end{itemize}
The unifying idea is that a mistake in kernel C does not produce a segfault in your program, because
there is no ``your program''. It produces an oops, and often it produces one somewhere else, later,
in code you did not write.

\subsection{Coding style, and why this repository does not use the kernel's}
\label{c4:app:a9}

The Linux kernel has a documented coding style, in \file{Documentation/process/coding-style.rst},
and it is not a matter of taste: patches are rejected for violating it, and
\file{scripts/checkpatch.pl} exists to check.

It is tabs, eight columns wide, with braces in K\&R placement:

\begin{ccode}
static int qa_fifo_put(struct qa_fifo *fifo, const u8 *data, size_t len)
{
	size_t n = 0;

	if (!fifo || !data)
		return 0;

	while (n < len && fifo->level < fifo->capacity) {
		fifo->data[(fifo->head + fifo->level) % fifo->capacity] = data[n++];
		fifo->level++;
	}

	return n;
}
\end{ccode}

\noindent \textbf{This repository is formatted differently}, with four spaces and braces on their
own line:

\begin{ccode}
static int qa_fifo_put(struct qa_fifo* fifo, const u8* data, size_t len)
{
    size_t n = 0;

    if (!fifo || !data) { return 0; }

    while (n < len && fifo->level < fifo->capacity)
    {
        fifo->data[(fifo->head + fifo->level) % fifo->capacity] = data[n++];
        fifo->level++;
    }

    return n;
}
\end{ccode}

\noindent That is the QAcademy house style, shared with every other course in the series and
enforced by \file{.clang-format}. The trade is deliberate and it is worth being honest about both
halves of it.

\textbf{What it buys} is that a reader moving between the Rust, C++ and Linux courses sees one
style, and that \code{make format} does the same thing everywhere.

\textbf{What it costs} is that none of the code in this repository would be accepted upstream, and
that a reader who goes from here to \file{drivers/} will find every file looks different. That
second cost is real, which is why this section exists, and why the two versions above are printed
side by side. The difference is entirely mechanical: eight-column tabs against four spaces, brace
placement, and \code{char *p} against \code{char* p}.

If you write a driver you intend to submit, run \code{scripts/checkpatch.pl --file} on it and follow
what it says. Nothing in this course's material depends on the formatting, and everything in it
would be correct in either style.
